\chapter{Tests und Validierung}
\label{ch:tests}

Das Projekt enthielt im aktuellen Stand zwei Testdateien. Sie pruefen die Konfiguration und das Verhalten der Geschwindigkeitsglättung sowie die Integration mit \texttt{GPSTrack} und dem Simulator.

\section{Ausgefuehrter Testbefehl}

Die Tests wurden mit dem folgenden Befehl ausgefuehrt:

\begin{lstlisting}
python -m unittest discover -s tests -v
\end{lstlisting}

\section{Ergebnis}

Die Test-Suite umfasst 14 Einzeltests. Alle 14 Tests sind erfolgreich durchgelaufen. Damit sind die Konfigurationsladung, die Aktivierung und Deaktivierung der Glättung, konstante Geschwindigkeit, Ausreisser, kurze Intervalle, grosse Zeitluecken, unregelmaessige Zeitabstaende sowie die Kompatibilitaet mit dem Simulator abgedeckt.

\begin{table}[htbp]
\centering
\caption{Abgedeckte Testbereiche}
\label{tab:testbereiche}
\begin{tabularx}{\textwidth}{@{}lX@{}}
\toprule
Testbereich & Gepruefte Aussage \\
\midrule
Konfigurationsladen & YAML-Werte werden eingelesen, Defaults greifen und ungueltige Werte fuehren zu Fehlern. \\
Aktivierte und deaktivierte Glättung & Die aktiven Spalten folgen entweder den Roh- oder den geglaetteten Werten. \\
Konstante Geschwindigkeit & Die Beschleunigung bleibt nahe null. \\
Einzelner Ausreisser & Ein grosser Geschwindigkeitssprung wird als Filterstuetzstelle ausgeschlossen. \\
Kurzes Intervall & Sehr kleine Zeitabstaende werden nicht als gueltige Stuetzstellen verwendet. \\
Grosse Zeitluecke & Neue Segmente beginnen an den richtigen Stellen, ohne ueber die Luecke zu interpolieren. \\
Unregelmaessige Intervalle & Der zeitabhaengige exponentielle Faktor reagiert auf unterschiedliche Zeitabstaende. \\
Simulator-Kompatibilitaet & \texttt{EBikeSimulator} verarbeitet das Track-DataFrame mit den aktiven Spalten korrekt. \\
\bottomrule
\end{tabularx}
\end{table}

\section{Interpretation}

Die Tests pruefen nicht die gesamte numerische Genauigkeit des Projekts, sondern die erwarteten Verhaltensweisen an den kritischen Schnittstellen. Insbesondere sichern sie ab, dass die Glättung keine pauschale Begrenzung einbaut und dass der Simulator die vom Track bereitgestellten aktiven Spalten ohne Zusatzannahmen weiterverarbeiten kann.
